\chapter*{Bibliographie / Webographie}
\phantomsection
\addcontentsline{toc}{chapter}{Bibliographie / Webographie}

Cette section regroupe les principales sources utilisées pour la compréhension du contexte, la rédaction du rapport et la justification des choix techniques. Les références internes concernent les informations fournies dans le cadre du stage et ne sont pas accessibles publiquement. Les références web correspondent principalement aux documentations officielles des technologies utilisées dans la solution \textbf{SomaScan}.

\section*{Bibliographie}
\phantomsection
\addcontentsline{toc}{section}{Bibliographie}

\begin{enumerate}
    \item SOMASTEEL. \textit{Informations internes relatives au contexte de l'entreprise, au processus logistique et au besoin de suivi des camions loués}. Documentation interne communiquée durant le stage, 2026.

    \item SOMASTEEL -- Direction des Systèmes d'Information. \textit{Échanges de cadrage, réunions de suivi et validation fonctionnelle du projet SomaScan}. Notes internes de stage, 2026.

    \item Projet SomaScan. \textit{Spécifications techniques internes : schéma de base de données, schéma API, architecture applicative et organisation des services}. Documents de conception et d'implémentation du projet, 2026.

    \item Mittelbach, F., Goossens, M., Braams, J., Carlisle, D. et Rowley, C. \textit{The LaTeX Companion}. Addison-Wesley, 2e édition, 2004.
\end{enumerate}

\section*{Webographie}
\phantomsection
\addcontentsline{toc}{section}{Webographie}

\small
\begin{longtable}{|p{4cm}|p{7.1cm}|p{3.2cm}|}
\hline
\textbf{Source} & \textbf{Utilisation dans le projet} & \textbf{Lien} \\
\hline
\endhead

Documentation Laravel & Référence pour le backend Laravel, l'organisation du projet, les routes, les contrôleurs, les migrations et le déploiement. & \url{https://laravel.com/docs/12.x} \\
\hline

Laravel Sanctum & Référence pour l'authentification par jetons utilisée par les applications mobile et web. & \url{https://laravel.com/docs/12.x/sanctum} \\
\hline

React Documentation & Référence pour la construction de l'application web d'administration avec des composants React. & \url{https://react.dev/learn} \\
\hline

Vite Documentation & Référence pour l'environnement de développement et le build de l'application React. & \url{https://vite.dev/guide/} \\
\hline

React Router Documentation & Référence pour le routage côté client et les modes de navigation utilisés dans l'application web. & \url{https://reactrouter.com/home} \\
\hline

Axios Documentation & Référence pour les appels HTTP entre l'application web et les API REST du backend. & \url{https://axios-http.com/docs/intro} \\
\hline

Tailwind CSS Documentation & Référence pour le système de classes utilitaires utilisé dans l'interface web. & \url{https://tailwindcss.com/docs/installation/using-vite} \\
\hline

Flutter Documentation & Référence pour le développement de l'application mobile destinée aux opérateurs terrain. & \url{https://docs.flutter.dev/} \\
\hline

Riverpod Documentation & Référence pour la gestion d'état et l'injection de dépendances dans l'application mobile Flutter. & \url{https://riverpod.dev/docs/introduction/getting_started} \\
\hline

Dio Package & Référence pour la communication HTTP côté mobile Flutter. & \url{https://pub.dev/packages/dio} \\
\hline

Flutter Secure Storage Package & Référence pour le stockage sécurisé des jetons d'authentification dans l'application mobile. & \url{https://pub.dev/packages/flutter_secure_storage} \\
\hline

Mobile Scanner Package & Référence pour la lecture des QR codes depuis l'application mobile. & \url{https://pub.dev/packages/mobile_scanner} \\
\hline

MySQL Documentation & Référence pour les principes généraux de base de données relationnelle utilisés avec Laravel. & \url{https://dev.mysql.com/doc/} \\
\hline

MariaDB Documentation & Référence complémentaire pour l'environnement de base de données compatible MySQL. & \url{https://mariadb.com/docs/} \\
\hline

GitHub Actions Documentation & Référence pour les workflows CI/CD et l'automatisation du déploiement. & \url{https://docs.github.com/en/actions} \\
\hline

GitHub Actions -- Writing workflows & Référence pour la structure des fichiers de workflow et le déclenchement des jobs. & \url{https://docs.github.com/en/actions/how-tos/write-workflows} \\
\hline

cPanel Documentation & Référence générale pour l'administration d'un hébergement web via cPanel. & \url{https://docs.cpanel.net/} \\
\hline

node-qrcode & Référence pour la génération de QR codes côté application web. & \url{https://github.com/soldair/node-qrcode} \\
\hline

jsPDF Documentation & Référence pour la génération de documents PDF contenant les QR codes imprimables. & \url{https://artskydj.github.io/jsPDF/docs/} \\
\hline

\caption{Sources web consultées}
\label{tab:webographie}
\end{longtable}
\normalsize

\vspace{0.3cm}
\noindent Toutes les sources web ont été consultées durant la rédaction du rapport.
